\documentclass{article}

\textwidth=6.5in
\textheight=8.5in
\voffset=-54pt
\oddsidemargin=0in
\evensidemargin=0in
\setlength{\parskip}{8pt}
\setlength{\parindent}{0pt}

\usepackage{amsmath, amssymb}
\usepackage{ifthen}

\usepackage[inline]{enumitem}
\setlist{topsep=0pt, parsep=8pt}

\usepackage[rgb,table]{xcolor}\convertcolorsUfalse
\usepackage{graphicx}

% change hyperref options
\PassOptionsToPackage{hyphens}{url}
\usepackage[bookmarks,colorlinks,linkcolor=black,citecolor=blue,pdfstartview=FitH,urlcolor=blue]{hyperref}
\hypersetup{pdfpagemode=UseNone}
\usepackage{bookmark}

\usepackage{graphicx}
\usepackage{multicol}
\usepackage{framed}
\usepackage{array}

\usepackage{icomma}

\usepackage{pifont}

\usepackage{marginnote}
\usepackage[colorinlistoftodos]{todonotes}
\renewcommand{\marginpar}{\marginnote}

\usepackage{soul}

\usepackage{pgfplots}
\pgfplotsset{compat=1.18}  
\pgfplotscreateplotcyclelist{mycolor}{black, blue, red, green!70!black, magenta, purple, cyan, orange }  
\pgfplotsset{
  disabledatascaling,
  cycle list=tikzcolor, 
  axis lines=middle,
  every axis plot/.style={no marks, thick},
  every axis/.style={
    enlarge x limits=false,
    enlarge y limits=auto,
    xlabel style={at={(current axis.right of origin)},anchor=west},
    ylabel style={at={(current axis.above origin)},anchor=south},
    % extra x and y ticks for asymptotes
    extra x tick style={grid=major, grid style={dashed,black}},
    extra y tick style={grid=major, grid style={dashed,black}},
    /pgf/number format/1000 sep={},
  },    
  tick label style = {font=\footnotesize, fill=\currentbgcolor, inner sep=0pt},    
  xticklabel shift=1pt,    
  scaled ticks=false,
  scaled y ticks=false,
  unbounded coords=jump,
  samples=101,
  minor tick length=0,
  scatter/classes={
    open={mark=*,draw=black,fill=white, mark size=2, thin},
    closed={mark=*,draw=black,fill=black, mark size=2, thin}
  },
  width=3in,
}
\pgfplotsset{open/.style={only marks, mark=*, fill=white, thin, mark size=2}}
\pgfplotsset{closed/.style={only marks, mark=*, draw=black, fill=black,
    thin, mark size=2}}
\pgfplotsset{labeled/.style={nodes near coords, only marks, mark=*, draw=black, fill=black, thin, mark size=2, point meta=explicit symbolic}}

\usepackage{tikz}
\usetikzlibrary{
  matrix,%
  % enables the snake paths
  decorations.pathmorphing,%
  % enables bigger arrows
  decorations.markings,%
  decorations.pathreplacing,%
  decorations.text,%
  calc,%
  patterns,%
  shapes.geometric,%
  arrows,
  decorations.shapes,
  positioning,
  plotmarks,
  intersections
}
\tikzset{>=latex} % sets default arrow type in tikz
\tikzset{font=\normalsize}
\tikzset{mark size=1.5pt, mark options=thin}
\tikzset{pin distance=4pt,
  every pin edge/.style={<-, thin, shorten <= -2pt}}

\interfootnotelinepenalty=10000
\renewcommand{\thefootnote}{(\arabic{footnote})}

\usepackage{multirow, bigdelim}

% change to \usepackage[sols]{handout} to see solutions
\usepackage[sols, workshop, grading]{handout}

\providecommand{\check}{}
\renewcommand{\check}{\ifmmode{\text{ \ding{51} }}\else{\ding{51} }\fi}

\newcommand\iscurrentenumi[1]{%
  \ifthenelse{\equal{#1}{\theenumi}}{}{#1}}
\setenumerate[1]{ref=\#\arabic*}
\setenumerate[2]{ref=\protect\iscurrentenumi{\theenumi}(\alph*)}
\newcommand\enumiiprefix[2]{%
  \ifthenelse{{\equal{#1}{\theenumi}}\AND{\equal{#2}{(\alph{enumii})}}}{}{}%
  \ifthenelse{{\equal{#1}{\theenumi}}\AND{\NOT{\equal{#2}{(\alph{enumii})}}}}{#2}{}%
  \ifthenelse{\equal{#1}{\theenumi}}{}{#1#2}%
}
\setenumerate[3]{ref=\protect\enumiiprefix{\theenumi}{(\alph{enumii})}\roman*}

\definecolor{purple}{rgb}{0.5,0,1.0}
\definecolor{hotpink}{rgb}{1, 0, 0.5}

\definecolor{skyblue}{rgb}{0, 0.5, 1}

\newcommand{\doedfinity}[2][\newpage]{
  \begin{problem}[#1]
    Do the Edfinity assignment ``Problem Set #2''.

    We encourage you to write down any scratch work here, as well as
    things you want your future self to remember. Nothing you write
    here will be graded; it's just for you!
  \end{problem}
}

\newcommand{\psrnewpage}{\newpage}
\newcommand{\psreflection}[2][]{

  \ifthenelse{\not\equal{#1}{nocite}}{
    \begin{enumerate}[resume]
    \item
      \begin{problem}[1in]
        Please cite any help you got on this problem set:
        \begin{itemize}[nosep]
        \item If you worked with other students, please list their names here.
        \item If you used resources other than those on our Canvas site, please list them here.
        \end{itemize}
      \end{problem}

    \end{enumerate}
  }

  \probonly{%
    #2

    \psrnewpage

    \emph{Reflection space.} It's a good habit to take a few minutes at the end of each pset to reflect on what you learned and jot down notes to your future self. Here are a few reflection questions to get you started. Nothing you write here will be graded; it's just for you!
    \begin{itemize}
    \item Take a few minutes to look back at the problems. How could you explain to someone else how to do each problem? What was your strategy, and how did you know to use that strategy? If there are problems where you don't feel able to answer this, write down which ones they are. Come back to them in a few days when we post the homework solutions!

      \vspace{1.5in}

    \item Take a look back at the learning objectives on today's worksheet; how are you feeling about each one? If there are ones you know you need more practice with, write those down here.

      \vspace{1.5in}

    \item What did you find most challenging about this material? Do you have any lingering questions about it? If so, write them down here so you don't forget them. At some point, stop by office hours or MQC to ask!

      \vfill
    \end{itemize}      
  }
}

\usepackage{fancyhdr}
\lhead{\textcolor{white}{This content is protected and may not be shared, uploaded, or distributed.}}
\rhead{}
\rfoot{\vspace*{\fill}\footnotesize\textcopyright{} \emph{Harvard Math Ma}}
\renewcommand{\headrulewidth}{0pt}
\pagestyle{fancy}
\begin{document}

\begin{center}
  \large\textsc{Problem Set 10 - The Derivative Function}
\end{center}

\begin{enumerate}
\item  \note{
    \begin{itemize}[nosep]
    \item On a graph, identify points where the function is non-differentiable (just corners for now).
    \item Interpret a derivative in context
    \end{itemize}
}

  \doedfinity{10}

\item  

  Let $f(x)=\dfrac{1}{x}$.
  \begin{enumerate}

  \item \label{derivative-of-1/x-graph-1/x}

    \begin{grading}
      1 point
    \end{grading}

    \begin{problem}[\vfill]
      Sketch the graph of $f(x)$. 
    \end{problem}

    \begin{solution}       
      \begin{center}
        \begin{tikzpicture}
          \begin{axis}[xtick=\empty, ytick=\empty, ymin=-3, ymax=3]
            \addplot+[domain=-3:3] {1/x};
          \end{axis}
        \end{tikzpicture}
      \end{center}
    \end{solution}

  \item \label{derivative-of-1/x-graph}

    \begin{grading}
      3 points: 1 for having the graph be negative everywhere, 1 for the graph being even, 1 for the graph increasing on $(0, \infty)$ and being very negative near 0
    \end{grading} 

    \begin{problem}[\vfill]
      Using only the graph you sketched in \ref{derivative-of-1/x-graph-1/x}, sketch a rough graph of $f'(x)$. 
    \end{problem}

    \begin{solution}       
      \begin{center}
        \begin{tikzpicture}
          \begin{axis}[xtick=\empty, ytick=\empty, ymin=-3, ymax=3]
            \addplot+[domain=-3:3, restrict y to domain=-5:0]
            {-1/x^2};
          \end{axis}
        \end{tikzpicture}
      \end{center}
      At any point on the graph of $f(x)$, the slope of the tangent line is negative, so $f'(x)$ is always negative.

      Because $f$ is odd, the slope of the tangent line at $x = a$ is equal to the slope of the tangent line at $x = -a$. That makes $f'$ even.

      As $x$ gets close to 0, the slope of the tangent line becomes very steep, so $f'$ is very negative when $x$ is close to 0.
    \end{solution}

  \item \label{derivative-of-1/x-calculate}

    \begin{grading} 
      4 points:
      \begin{itemize}[nosep]
      \item 3 for correctly simplifying the difference quotient $\dfrac{ f(x+h) - f(x) }{h}$ (give partial credit if the simplification is partially correct)
      \item 0.5 for using limit notation correctly (writing $\displaystyle \lim_{h \to 0}$ on each line until they actually take the limit)
      \item 0.5 for correct limit
      \end{itemize}
    \end{grading} 

    \begin{problem}[\nstretch{2.25}]
      Use the limit definition of derivative to find a formula for $f'(x)$.
    \end{problem}

    \begin{solution}  
      \begin{align*}
        f'(x) &= \displaystyle \lim_{h \rightarrow 0} \frac{ f(x+h) - f(x) }{h} \\
        &= \displaystyle \lim_{h \rightarrow 0} \frac{ \frac{1}{x+h} - \frac{1}{x} } {h} \\
        &= \displaystyle \lim_{h \rightarrow 0} \frac{ \frac{ x - (x+h)}{x(x+h)}}{h}\\
        &= \displaystyle \lim_{h \rightarrow 0}  \frac{-h}{x(x+h)} \cdot \frac{ 1 } {h} \\
        &= \displaystyle \lim_{h \rightarrow 0}\frac{ -1}{x(x+h)} \\
        & =\boxed{ -\frac{1}{x^2}}\\
      \end{align*}
    \end{solution}

  \item

    \begin{grading}
      Nothing to grade
    \end{grading}

    \begin{problem}
      \check Check that your formula for $f'(x)$ is consistent with the graph you drew in \ref{derivative-of-1/x-graph}. (You don't need to write anything.)
    \end{problem}

  \item
    \begin{grading} 
      3 points: 1 for solving correctly, 1 for labeling graph of $f$ appropriately, 1 for labeling graph of $f'$ appropriately 
    \end{grading} 

    \begin{problem}[\stretch{0.5}]
      Find all values of $x$ for which $f'(x) = -2$. On the graphs you sketched in \ref{derivative-of-1/x-graph-1/x} and \ref{derivative-of-1/x-graph}, label these points, as well as where we visualize $-2$ on each graph.
    \end{problem}

    \begin{solution}
      Since we found that $f'(x) = - \dfrac{1}{x^2}$, we just want to solve
      \begin{align*}
        - \dfrac{1}{x^2} &= - 2 \\
        \intertext{Negating and taking the reciprocal of both sides,}
        x^2 &= \frac{1}{2} \\
        x &= \boxed{\pm \frac{1}{\sqrt{2}}}
      \end{align*}
      On the graph of $f$, this describes the points where the tangent line has slope $-2$:
      \begin{center}
        \begin{tikzpicture}
          \begin{axis}[xtick={-0.707, 0.707},
            xticklabels={$-\frac{1}{\sqrt{2}}$, $\frac{1}{\sqrt{2}}$},
            ytick=\empty, ymin=-3, ymax=3, clip=false]
            \addplot+[domain=-3:3, restrict y to domain=-3:3, samples=109] {1/x} node[right]{$y = f(x)$}; %
            \addplot+[closed, skyblue] coordinates { ({-sqrt(1/2)},
              {-sqrt(2)}) ({sqrt(1/2)}, {sqrt(2)}) }; %
            \addplot+[domain={-sqrt(1/2)-0.5}:{-sqrt(1/2)+0.5},
            hotpink] {-2*(x+sqrt(1/2))-sqrt(2)} node[right,
            fill=white]{slope $-2$}; %
            \addplot+[domain={sqrt(1/2)-0.5}:{sqrt(1/2)+0.5}, hotpink]
            {-2*(x-sqrt(1/2))+sqrt(2)} node[pos=0, anchor=south east, fill=white]{slope
              $-2$}; %
          \end{axis}
        \end{tikzpicture}
      \end{center}
      On the graph of $f'$, this describes where the value of $f'$ (which is the value on the vertical axis) is $-2$:
      \begin{center}
        \begin{tikzpicture}
          \begin{axis}[xtick={-0.707, 0.707},
            xticklabels={$-\frac{1}{\sqrt{2}}$, $\frac{1}{\sqrt{2}}$},
            ytick={-2}, ymin=-3, ymax=3, clip=false]
            \addplot+[domain=-3:-0.57735] {-1/x^2}; %
            \addplot+[domain=0.57735:3]
            {-1/x^2} node[right]{$y = f'(x)$}; %
            \addplot+[closed, skyblue] coordinates { ({-sqrt(1/2)},
              -2) ({sqrt(1/2)}, -2) }; %
          \end{axis}
        \end{tikzpicture}
      \end{center}
    \end{solution} 

  \end{enumerate}

\item  

  Let $a$, $b$, and $c$ be arbitrary constants. 

  \begin{enumerate}

  \item \label{derivative-of-general-quadratic}

    \begin{grading}
      4 points, same breakdown as \ref{derivative-of-1/x-calculate}
    \end{grading} 

    \begin{problem}[\vfill\newpage]
      Use the limit definition of derivative to find the derivative of
      $g(x) = {a}x^2 + bx + c$. 
    \end{problem}

    \begin{solution}       
      \begin{align*} 
        g'(x) &= \displaystyle \lim_{h \rightarrow 0} \frac{ g(x+h) - g(x)}{h} \\
        &= \displaystyle \lim_{h \rightarrow 0} \frac{ a(x+h)^2+b(x+h) +c - (ax^2+bx+c) } {h}\\
        & = \displaystyle \lim_{h \rightarrow 0} \frac{ ax^2+2axh+h^2 +bx+bh+c -ax^2-bx-c}{h}\\
        & = \displaystyle \lim_{h \rightarrow 0} \frac{ 2axh+bh+ h^2}{h} \\
        &= \displaystyle \lim_{h \rightarrow 0} \frac{h(2ax+b+h)}{h}\\
        &= \displaystyle \lim_{h \rightarrow 0} (2ax +b+h) \\
        &= 2ax+b
      \end{align*} 
    \end{solution} 

  \item \label{derivative-of-constant}

    \begin{grading}
      The remaining parts of this problem are 1 point each. 
    \end{grading}

    \begin{problem}[\vfill]
      Explain what your answer to \ref{derivative-of-general-quadratic} tells you about the derivative of a constant function $g(x) = c$. (What values of $a$ and $b$ does this correspond to?)
    \end{problem}

    \begin{solution}
      Starting with the function $g(x)=ax^2+bx+c$, we get a constant function $g(x)=c$ if $a$ and $b$ are both zero.
      Based on the formula we found in part \ref{derivative-of-general-quadratic}, we can deduce that
      the derivative of a constant function is $g'(x)=0$.
    \end{solution}

  \item

    \begin{problem}[\vfill]
      Explain in words why your answer to \ref{derivative-of-constant} makes sense graphically.
    \end{problem}

    \begin{solution}
      This makes sense because the graph of a constant function is a horizontal line, and so the slope of any
      tangent line will always be zero.
    \end{solution}

  \item \label{derivative-of-linear}

    \begin{problem}[\vfill]
      Explain what your answer to \ref{derivative-of-general-quadratic} tells you about the derivative of a linear function $g(x) = bx + c$.%
    \end{problem}

    \begin{solution}
      The function $g(x)=bx+c$ corresponds to the case when $a=0$, and based on the formula we found in part 
      \ref{derivative-of-general-quadratic}, the derivative would be $g'(x)=b$. 
    \end{solution}

  \item

    \begin{problem}[\vfill\newpage]
      Explain in words why your answer to \ref{derivative-of-linear} makes sense graphically.
    \end{problem}

    \begin{solution}
      This makes sense because the graph of a linear function is a straight line, and so the slope of any tangent line
      will always be the same number.
    \end{solution}

  \end{enumerate}  

\item  

  Let $j(x) = 5 - (x - 3)^2$.

  \begin{enumerate}
  \item \label{shifted-quadratic-graph}

    \begin{grading}
      2 points: 1 for a downward parabola with vertex at $(3, 5)$, 1 for a correct explanation, 1 for intercepts in roughly the right places (they'll find these in \ref{shifted-quadratic-intercepts})
    \end{grading}

    \begin{problem}[\vfill]
      Use what you know about shifts and stretches to sketch the graph of $j(x)$, and explain how it's related to the graph $y = x^2$. Please label any significant values on the axes. 
    \end{problem}

    \begin{solution}
      First, we can rewrite $j(x)$ as $j(x)=-(x-3)^2 + 5$. Now it is clearer that the graph of $j$ is the graph of $y=x^2$ but
      shifted to the right by 3 units, reflected over the $x$-axis, and shifted up by 5 units.

      \begin{tikzpicture}
        \begin{axis}[xtick={0.764, 3, 5.236}, ytick={-4,5}, xticklabels={$3-\sqrt{5}$, $3$, $3+\sqrt{5}$}]
          \addplot[domain=0:6] {-(x-3)^2+5};
        \end{axis}
      \end{tikzpicture}
    \end{solution}

  \item \label{shifted-quadratic-intercepts}

    \begin{grading}
      2 points: 0.5 for the $y$-intercept, 1.5 for the $x$-intercepts
    \end{grading}

    \begin{problem}[\vfill]
      Find the $x$- and $y$-intercepts of $y = j(x)$. If your graph in \ref{shifted-quadratic-graph}  doesn't reflect these accurately, go back and adjust your graph.
    \end{problem}

    \begin{solution}
      To find the $y$-intercept, we set $x=0$ and solve for $y$: $y= 5-(0-3)^2 = -4$. To find the $x$-intercept(s),
      we set $y=0$ and solve for $x$:
      \begin{align*}
        0 &= 5-(x-3)^2 \\
        (x-3)^2 &= 5 \\
        x-3 &= \pm \sqrt{5} \\
        x &= 3 \pm \sqrt{5}
      \end{align*}

    \end{solution}

  \item \label{shifted-quadratic-positive-derivative}

    \begin{grading}
      The remaining parts of this problem are 1 point each.
    \end{grading}

    \begin{problem}[\nstretch{0.5}]
      Based on the graph you sketched, for what values of $x$ is $j'(x)$ positive?
    \end{problem}

    \begin{solution}
      Based on the graph, $j'(x)$ is positive when $x<3$.
    \end{solution}

  \item

    \begin{problem}[\stretch{0.5}]
      Use algebra to rewrite $j(x)$ in the form $ax^2 + bx + c$.
    \end{problem}

    \begin{solution}
      \begin{align*}
        j(x) &= 5-(x-3)^2 \\
        &= 5-(x^2 - 6x + 9) \\
        &= 5 - x^2 + 6x - 9 \\
        &= -x^2 + 6x - 4
      \end{align*}
    \end{solution}

  \item \label{shifted-quadratic-j'}

    \begin{problem}[\stretch{0.25}]
      Use the result of \ref{derivative-of-general-quadratic} to find $j'(x)$. 
    \end{problem}

    \begin{solution}
      We know that if $f(x)=ax^2+bx+c$, then $f'(x)=2ax+b$. Since $j(x) = -x^2 + 6x - 4$, where
      $a = -1$, $b=6$, and $c=-4$, we know
      $j'(x) = 2(-1)x + 6 = -2x+6$. 
    \end{solution}

  \item \label{shifted-quadratic-j'(2)}

    \begin{problem}[\stretch{0.25}]
      Use your answer to \ref{shifted-quadratic-j'} to find $j'(2)$. 
    \end{problem}

    \begin{solution}
      $j'(2) = -2(2)+6 = 2$. 
    \end{solution}

  \item

    \begin{grading}
      There are several things they could mention:
      \begin{itemize}[nosep]
      \item that their formula for $j'(x)$ agrees with where they said $j'(x)$ is positive
      \item that their sign of $j'(2)$ agrees with where they said $j'(x)$ is positive
      \item that their sign of $j'(2)$ agrees with their sketch of $j(x)$
      \end{itemize}
      Give full credit if they mention any one of these.
    \end{grading}

    \begin{problem}[\stretch{0.5}]
      \check Are your answers to this problem consistent with each other? Explain.
    \end{problem}

    \begin{solution}
      Yes. In part (c), we determined that $j'(x)$ is positive when $x<3$, and this agrees with the formula
      we found, $j'(x) = -2x+6$, as well as the value we calculated, $j'(2) = 2$. 
    \end{solution}

  \end{enumerate}

\item 
  \begin{minipage}[t]{1.0\linewidth-2in}
    \setlength{\parskip}{8pt}
    Marvin is standing at one corner of a giant rectangular field. He'd like to get to the opposite corner, which is 400 meters east and 600 meters north of him. The southern third of the field is covered with snow; the rest is just grass. Marvin walks through snow at a rate of 0.5 m/s and walks on grass at a rate of 1.5 m/s. Rather than taking a straight-line path to the opposite corner, Marvin decides to take a straight path through the snow to a point $x$ meters east and 200 meters north of where he started, and then take a straight path from there to his destination. 
  \end{minipage}
  \hfill
  \parbox[t]{1.5in}{
    \vspace{0pt}
    \begin{tikzpicture}[x=3.5cm, y=2cm]
      \draw (0, 0) -- (1, 0) -- (1, 3) -- (0, 3) -- cycle;
      \draw[dashed] (0, 1) -- (1, 1) node[anchor=south east]{\emph{grass}}
      node[anchor=north east]{\emph{snow}};
      \draw[very thick] (0, 0) -- (0.2, 1) -- (1, 3);
      \draw[decorate, decoration={brace}] (0, 1) --
      node[midway, above]{$x$} (0.2, 1);   
    \end{tikzpicture}
  }    

  \probonly{\emph{Continued on next page}}

  \pspace{\newpage}

  \begin{enumerate}
  \item
    \begin{grading}
      3.5 points:
      \begin{itemize}[nosep]
      \item 1 for labeling the picture
      \item 0.5 for stating the goal
      \item 1 for relating the goal to distance and rate
      \item 1 for expressing the relevant distance in terms of $x$
      \end{itemize}
    \end{grading}

    \begin{problem}[\vfill]
      Express the time (in seconds) that it takes Marvin to walk through the snow as a function of $x$. 

      (Remember the guidelines for presenting your work in such problems. If you've forgotten, look back at \href{https://people.math.harvard.edu/~jjchen/mathma/docs/homework/PS06.html}{Problem Set 6}.)
    \end{problem}

    \begin{solution}
      The \hl{\textbf{goal}} is to express the time it takes Marvin to walk through the snow as a function
      of $x$, the length indicated on the diagram.

      We \hl{\textbf{know}}:
      \begin{itemize}
        \item The point where Marvin exits the snow is 200 meters north of where he starts.
        \item Marvin's walking speed in the snow is 0.5 m/s.
      \end{itemize}

      Let's \hl{\textbf{draw a picture and define variables.}}
      \begin{center}
        \begin{tikzpicture}[x=3.5cm, y=2cm]
          \draw (0, 0) -- (1, 0) -- (1, 1) -- (0, 1) -- cycle; 
          \draw[very thick] (0, 0) -- node[midway, right, red] {$d_1$} (0.2, 1); 
          \draw[decorate, decoration={brace}, red] (0, 1) -- node[midway, above]{$x$} (0.2, 1); 
          \draw[|-|, xshift=-1.5em, red] (0, 0) -- node[midway, fill=white]{$200$} (0, 1); 
        \end{tikzpicture}
      \end{center}
      Here, $x$ is the length described in the problem and $d_1$ is the length that Marvin walks through the snow.
      Let's call the time it takes Marvin to walk through the snow $t_1$. 

      Now let's \hl{\textbf{write a formula}} for $t_1$. Since
      $\text{distance} = \text{speed} \times \text{time}$, $t_1 = \frac{d_1}{0.5}$.

      To express $t_1$ as a function of $x$, we need to \hl{\textbf{relate}} $d_1$ to $x$. We can do this by applying
      the Pythagoren theorem to the right triangle in the picture, so 
      \begin{align*}
        d_1^2 &= x^2 + 200^2 \\
        d_1 &= \sqrt{x^2 + 200^2}.
      \end{align*}
      Substituting this expression for $d_1$ into our formula for $t_1$, we get that 
      \fbox{$t_1 = \dfrac{\sqrt{x^2 + 200^2}}{0.5}$.}
    \end{solution}

  \item

    \begin{grading}
      2.5 points:
      \begin{itemize}[nosep]
      \item 0.5 for stating the goal
      \item 0.5 for recognizing that they can add time through the grass and time through the snow
      \item 0.5 for relating the time through the snow to distance and rate
      \item 1 for expressing the relevant distance in terms of $x$
      \end{itemize}
    \end{grading}

    \begin{problem}[\vfill\newpage]
      Express the total time (in seconds) it takes Marvin to get from his starting corner to the opposite corner as a function of $x$.
    \end{problem}

    \begin{solution} 
    The \hl{\textbf{goal}} is to express the total time Marvin takes to go from his starting position 
    to the opposite corner as a function of $x$. 
    
    In addition to what we wrote down in part (a), we also \hl{\textbf{know}}:
    \begin{itemize}
      \item The opposite corner is 400 meters east and 600 meters north of the starting position
      \item Marvin's walking speed on grass is 1.5 m/s.
    \end{itemize}

    Let's \hl{\textbf{draw a picture and define variables.}}
    \begin{center}
      \begin{tikzpicture}[x=3.5cm, y=2cm]
        \draw (0, 0) -- (1, 0) -- (1, 3) -- (0, 3) -- cycle; 
        \draw[dashed] (0, 1) -- (1, 1); 
        \draw[very thick] (0, 0) -- (0.2, 1) node[midway, right, red] {$d_1$} -- node[midway, right, red] {$d_2$} (1, 3); 
        \draw[decorate, decoration={brace}, red] (0, 1) -- node[midway, above]{$x$} (0.2, 1); 
        \draw[decorate, decoration={brace}, red] (0.2, 1) -- node[midway, above]{$400 -x$} (1, 1); 
        \draw[|-|, xshift=-1.5em, red] (0, 0) -- node[midway, fill=white]{$200$} (0, 1);
        \draw[|-|, xshift=1.5em, red] (1, 1) -- node[midway, fill=white]{$400$} (1, 3);
      \end{tikzpicture}
    \end{center}

    We define $x$ and $d_1$ the same as before, and we let $d_2$ be the distance that Marvin walks through the grass.
    We will use $t_1$ to represent the time it takes Marvin to walk through the snow, $t_2$ to represent the time it takes
    Marvin to walk through the grass, and $t_\text{total}$ to represent the total time it takes Marvin to get to the
    opposite corner.

    Now let's \hl{\textbf{write a formula}} for $t_\text{total}$. First, we can say that $t_\text{total} = t_1 + t_2$. 

    To express $t_\text{total}$ as a function of $x$, we need to \hl{\textbf{relate}} $t_1$ and $t_2$ to $x$. We already
    know that $t_1 = 2\sqrt{x^2 + 200^2}$. We need to find a way to write $t_2$ in terms of $x$.  Since
    $\text{distance} = \text{speed} \times \text{time}$, $t_2 = \frac{d_2}{1.5}$. We can relate $d_2$ to $x$
    by applying the Pythagorean theorem to a right triangle in our picture.
    \begin{align*}
      d_2^2 &= (400-x)^2 + 400^2 \\
      d_2 &= \sqrt{(400-x)^2 + 400^2}.
    \end{align*}
    This means that $t_2 = \dfrac{\sqrt{(400-x)^2 + 400^2}}{1.5}$, and so
    \[t_\text{total} = \dfrac{\sqrt{x^2 + 200^2}}{0.5} + \dfrac{\sqrt{(400-x)^2 + 400^2}}{1.5}.\]
    
    \end{solution}
  \end{enumerate}

\end{enumerate}
\renewcommand{\psrnewpage}{{\centering\rule{5in}{0.4pt}}}

\psreflection{For next class, read \S 5.3, 8.2.}
\end{document}
